\Artigo%
{miscelanea}%
{Erros, chapuzas e desatinos}%
{Ánxel Costas}%
{Historias mal feitas para facelo mellor.}

\begin{multicols}{2}

Non é mentira algunha que moitísima xente quede abraiada coa cantidade de  %ese algunha talvez sexa ningunha, sóame estraño, pero de momento deixémolo así
avances e descubrimentos que veñen da man da física, capaces de desafiar o
sentido común da xente menos entendida. E é que, ás veces, a física parece
maxia e nós, os (case) físicos, magos. Pero estes avances non se fan da noite á
mañá, pois hai unha gran cantidade de persoas que traballaron e traballan duro
durante moitos anos para acadar estes fitos na nosa historia. Mais non todos
estes esforzos foron parar a bo porto, máis ben ao contrario, foron ocultos no
fondo do caixón da vergoña no soto da humillación e tirada a chave polo pozo do
esquecemento. E creo que é hora de destapar os trapos sucios que levamos
agochando tanto tempo, pois a humildade é algo que podemos esquecer facilmente
ao crer que comprendemos o Universo e as súas forzas, e cómpre lembrala de
cando en vez. É por iso que veño contarvos a vós, meus ben queridos lectores,
varias «cafradas» que se fixeron en física, para que saibades que unha
integral mal feita ou un signo mal posto non é tan grave, a non ser que sexa
financiado, creado, exposto, explotado e posteriormente oculto. Todo o demais
son erratas que pode ter calquera persoa ao pensar rápido (todos sabemos que é
por iso e non porque as mates ás veces se nos fan costa arriba). Así que, %non estou seguro de que a abreviatura mates sexa correcta en galego, pero de momento deixémolo así
parodiando a frase coa que Lord Kelvin meteu a pata: «Está todo ben feito en
física, salvo un par de problemas que hai que ocultar»; vouvos contar tres
destes problemas que foron «ocultos» (ou non tanto), a modo de tres contos de
Nadal.

\subsection*{Os raios N}

O gran protagonista desta historia (que rolda o ridículo colectivo) foi
Prosper-René Blondlot, un dos oito físicos da Academia das Ciencias de Francia.
Mentres traballaba cos raios X en 1903, percibiu variacións do brillo cando
realizaba fotografías dun entreferro, o que interpretou como un novo tipo de
radiación, que decidiu chamar raios N (pola universidade de Nancy onde
traballaba). Este descubrimento levou a 300 artigos e 120 físicos afirmar a
existencia desta nova radiación, presente en case todos os corpos, salvo en
madeiras verdes e algúns metais tratados. Ben, parece que está todo correcto e
non hai ningún problema se tanta xente estaba de acordo. Entón, onde está a
trampa? Todo foi destapado por Robert W. Wood, quen xa tiña fama de desmentir
resultados excéntricos. Nunha visita ao laboratorio de Blondlot para a
demostración dos raios N, Wood (nun momento de lucidez propia dun pillabán),
aproveitándose da escuridade necesaria para o experimento, decidiu quitar un
prisma que amplificaba supostamente este fenómeno sen que ninguén se decatase.
O resto da sala dixo que seguía vendo estes raios N, o cal xa era unha alarma
xigante da mentira que era todo, pero para confirmar, tamén decidiu cambiar un
arquivador metálico por un andel de madeira que, supostamente, non debería
emitir esta radiación. As demais persoas presentes seguían vendo o efecto
luminoso, co cal xa quedou máis que claro que toda esa xente era uns farsantes
ou, como moito, unhas vítimas da autosuxestión (colectiva). Pero quen somos nós
para xulgar, se tamén temos «falseado» resultados nos laboratorios para
aprobar as materias dos profesores? Esta historia non é máis que unha lección de % cambiado "do" -> "dos"
que os resultados acadados non sempre son os que cremos que son, e que ter un % "de que" -> "que"?
compañeiro de laboratorio capaz de levarnos a contraria é necesario para pórnos
os pés na terra.

\subsection*{Neutrinos máis rápidos que a luz}

Imos ver se un século despois deprendemos algo do ridículo caso anterior. Esta
historia sitúase no 2011, sendo o CERN, e concretamente o experimento OPERA, o
centro do foco, os físicos que traballaban neste proxecto tiveron un resultado
inusual. Resulta que os neutrinos medidos semellaban superar a velocidade da
luz, 60 nanosegundos máis rápidos. Se isto fose medido por algún alumno da   %a velocidade, que eu saiba, non se mide en segundos
facultade de Física, o máis probable é que o devandito alumno fixese maxia e os
valores medidos cambiasen de golpe, dando lugar a uns resultados fermosísimos
dentro do esperado, pero no CERN traballa xente máis seria e querían saber a
que se debía tal anomalía. E é que os resultados presentáronos como iso, unha
anomalía, curáronse en saúde antes de anunciar un descubrimento (vese que a
historia anterior foi escarmento dabondo). Despois de varios meses investigando
o porqué destes resultados, atopáronse os culpables: un cable de fibra óptica % "deses" -> "destes"?, cambiado "atopouse" -> "atopáronse"
mal conectado e un oscilador de reloxo defectuoso, os cales introducían un erro
sistemático nos resultados. Ás veces a culpa non é dos físicos, senón dos
materiais, pero o dano xa estaba feito: miles de teorías conspiranoicas,
queixas populares sobre os fondos para os experimentos para un cable mal
conectado, moitísimos filmes de ciencia ficción de mala calidade... Quizais por
isto teñamos pánico a presentar uns resultados (públicos ou ao profesor) sen
saber se son correctos, pois os acertos traen avances, pero os desacertos traen
moitas dores de cabeza. A ensinanza desta historia é revisar a montaxe dos
experimentos se non queres pasar os próximos días pelexando cuns valores
carentes de sentido. Despois deste incidente, comezou a soar a frase «un cable
frouxo pode ser máis rápido que a luz».

\subsection*{A explosiva nova do Mars Climate Orbiter} % cambiado "Orber" -> "Orbiter"

Esta pequena anécdota remóntase a 1999. Aquí non hai nin autosuxestión nin
problemas coa montaxe do experimento, estamos ante unha situación de descoido
(por non dicir estupidez) puramente humana. Un pequeno erro que custou 301
millóns de euros (axustando á inflación de hoxe en día, uns 535 millóns),
tirados ao lixo, ou máis ben, á atmosfera de Marte. A misión da sonda espacial
era estudar a atmosfera do planeta vermello durante dous anos terrestres, pero
nunca puido realizarse. A medida que a sonda se acercaba ao planeta, os
controladores decatáronse de que había que realizar demasiadas correccións, máis % engadido (decatáronse) "de" (que)
das esperadas, e aí comezaron as sospeitas, polo que decidiron pedir unha
investigación. Os responsables do proxecto, como bos responsables, decidiron
desestimar a petición de investigación, e posteriormente evadiron as acusacións
alegando que a petición de investigación non se realizara formalmente
(comportamento máis propio dun político que de xente adicada á ciencia). Pero, %adicar: forma menos recomendábel por dedicar
cal foi o accidente? Pois que a sonda foise apartando da órbita esperada a % "foise": pronome ben colocado?, engadido acento gráfico en "órbita"
medida que se acercaba a Marte, até que pasou a só 57 km de altura, cando se
esperaba que pasase a 150 km. Por suposto, a sonda quedou completamente
esnaquizada pola fricción coa atmosfera. Agora que xa coñecemos o accidente,
como bos físicos deberiamos preguntarnos cal foi o problema. As unidades. As
malditas unidades. Esas mesmas que nos levan dicindo dende o colexio «47 que?
mazás? centímetros?». Un non é consciente do importantes que son até que perde
301 millóns en apenas segundos. A sonda empregaba unidades do Sistema
Internacional, pero os controladores da NASA traballaban con unidades imperiais
(se os estadounidenses non levasen a contraria no sistema de medida, este erro
non tería acontecido, ou polo menos non por este motivo). A diferenza entre
ambos os sistemas foise acrecentando a medida que a sonda avanzaba pola súa
órbita, a pesar das correccións dos controladores, até acabar estrelándose. Que
podemos aprender de todo isto? Como primeiro e mínimo, non fiarse da xente que
empregue unidades de medida distintas ao Sistema Internacional e, como segundo,
que hai que revisar as contas. Sempre revisar as contas, pois nunca sabes onde
podes atopar un erro que pode arruinar o teu exercicio (ou arruinarte
economicamente). % engadido ")"

Espero que estas tres pequenas historias fagan manternos humildes, e facernos
recapacitar sobre as vergoñas históricas que agochamos e das cales podemos
deprender moito. E lembrade, meus queridísimos lectores, que cando acabedes de
ler este pequeno artigo deberiades de ter presente as unidades, montar ben o
experimento e ter un bo compañeiro que vos corrixa de ser necesario, para
evitar ter erros, chapuzas e desatinos.

\end{multicols}
